\documentclass[11pt]{article}
\usepackage[T1]{fontenc}
\usepackage[utf8]{inputenc}
\usepackage{lmodern,amsmath,amssymb,amsthm,mathtools}
\usepackage[margin=27mm]{geometry}
\usepackage[hidelinks]{hyperref}
\newtheorem{theorem}{Theorem}
\newtheorem{proposition}[theorem]{Proposition}
\newcommand{\Z}{\mathbb Z}
\newcommand{\C}{\mathbb C}
\newcommand{\R}{\mathbb R}
\title{The EP817 residue source in the original supported quotient:
exact cochain maps, metric losses, and arithmetic transitions}
\author{The Clankers}
\date{19 September 2026}
\begin{document}
\maketitle

This note integrates the supplied \emph{Residue prices, constrained composition,
and nonperiodic optimal schedules} of 17 September 2026
\cite{continuation}.  Its full proof, original LaTeX, and exact certificates
are preserved without alteration in the accompanying source subtree.
The maps below implement the internal quotient and receiving-kernel construction
of the original Split-Zero support diagram, equations (D6)--(D8)
\cite{support}.  All word multiplicities, actual integer values, moduli,
support labels, and Euclidean word metrics are retained.  In the
composition results the coefficient arity is an integer \(q\geq2\);
the applications at \(q=3,5\) therefore have the original cut constant \(q-1\).

\section{The actual source and its two cochain complexes}

Let \(A\) be a fixed nonempty finite set of distinct positive integers and
let \(b>\sum_{a\in A}a\) be the original radix.  Set
\[
 \Omega=\{0,1,2\}^{A},\qquad
 v(c)=\sum_{a\in A}c_a a,\qquad
 \bar v(c)=v(c)\bmod b .
 \tag{E1}
\]
For every \(S\subseteq\Omega\) use the actual images
\[
 X_S=v(S)\subset\Z,\quad R_S=\bar v(S)\subset\Z/b\Z,\quad
 V_S=\Z[S],\quad E_S=\Z[X_S],\quad \bar E_S=\Z[R_S].
 \tag{E2}
\]
Here square brackets denote free abelian groups on the indicated distinct
bases, including a basis vector \(e_0\) whenever the image contains zero.
The maps on these displayed bases are
\[
 q_S(e_c)=e_{v(c)},\qquad
 \pi_S(e_x)=e_{x\bmod b},\qquad p_S=\pi_Sq_S.
 \tag{E3}
\]
All three have their stated codomains by (E2), and both \(q_S,p_S\) are
surjective.  Their kernels are
\[
 B_S=\ker q_S,\qquad B'_S=\ker p_S,\qquad B_S\subseteq B'_S .
 \tag{E4}
\]
For an empty image its free group is zero and the same formulas apply.

Define two complexes concentrated in degrees \(-1,0\):
\[
 C_S=[\,B_S\overset{\iota_S}{\longrightarrow}V_S\,],\qquad
 D_S=[\,B'_S\overset{\iota'_S}{\longrightarrow}V_S\,].
 \tag{E5}
\]
The map \(f_S:C_S\to D_S\) is inclusion in degree \(-1\) and identity
in degree zero.  It is a cochain map because
\(\iota'_S(B_S\hookrightarrow B'_S)=\iota_S\).

\begin{theorem}[The exact receiving kernel in the original complexes]
The maps in (E3) induce the isomorphisms and exact sequence
\[
 \begin{gathered}
 H^0(C_S)\overset{\sim}{\longrightarrow}E_S,\quad[c]\longmapsto q_Sc,
 \qquad
 H^0(D_S)\overset{\sim}{\longrightarrow}\bar E_S,\quad[c]\longmapsto p_Sc,\\
 H^{-1}(C_S)=H^{-1}(D_S)=0,\qquad H^0(f_S)=\pi_S,\\
 0\longrightarrow B_S\longrightarrow B'_S
       \overset{q_S}{\longrightarrow}\ker\pi_S\longrightarrow0,\qquad
 B'_S/B_S\overset{\sim}{\longrightarrow}\ker H^0(f_S),
       \quad[c]\longmapsto q_Sc .
 \end{gathered}
 \tag{E6}
\]
\end{theorem}
\begin{proof}
Injectivity of the two differentials gives the degree-\(-1\) statements.
Their degree-zero cohomologies are \(V_S/B_S\) and \(V_S/B'_S\).
The maps \(q_S,p_S\) are onto with exactly those kernels, proving the two
displayed isomorphisms.  Their commutative square is (E3), so their induced
degree-zero map is precisely \(\pi_S\).

If \(c\in B'_S\), then \(\pi_Sq_Sc=p_Sc=0\), so the third map of the
sequence has the asserted codomain.  Its kernel is exactly \(B_S\).
To prove its surjectivity without an unspecified lift, order the original
coefficient words lexicographically using the increasing order on \(A\).
For each \(x\in X_S\), let \(c_x\) be the first word in \(S\) with value
\(x\).  An element \(h=\sum_xh_xe_x\in\ker\pi_S\) is the image under
\(q_S\) of \(\sum_xh_xe_{c_x}\); that lift belongs to \(B'_S\) because
its \(p_S\)-image is \(\pi_Sh=0\).  This proves the exact sequence and
the final isomorphism.
\end{proof}

For \(S\subseteq T\), use basis inclusions in each of \(V,E,\bar E\).
The squares with \(q\) and \(\pi\) commute on every basis vector.
Consequently \(B_S\to B_T\) and \(B'_S\to B'_T\) are the restrictions
of \(V_S\to V_T\), and the complexes and (E6) form actual diagrams.
Use the join-semilattice
\[
 L=\{\bot\}\sqcup\mathcal P(\Omega),\qquad S\vee T=S\cup T,
 \tag{E7}
\]
with a distinct bottom \(\bot\) and zero fibre there.  The label
\(\varnothing\in\mathcal P(\Omega)\) remains distinct from \(\bot\),
despite its zero amplitude group.  Over each displayed integral diagram
the split carrier consists of \((S,x)\), with addition by transport to
\(S\cup T\), supported scalar multiplication \(r^\bullet(S,x)=(S,rx)\),
and \(\tau(S,x)=(\bot,0)\).  These are the original operations
of \cite[(D2)]{support}.  Commutativity of the above squares proves
that the reconstructed \(q,\pi,f\) are split-linear.

In particular the all-label zero-amplitude preimage of the reconstructed
\(\pi\) is
\[
 \coprod_{S\in L}\ker\pi_S .
 \tag{E8}
\]
The preimage of its global zero contains only the bottom fibre, since its
label map is identity.  The exact comparison is the inclusion of that
bottom fibre into (E8); (E6) specifies each nonbottom fibre of (E8).
A supported boundary \((S,c)\) goes to \((S,0)\) at the same label.
Also the residue-amplitude map
\(\bar E_S\to\Z/b\Z,\ \sum_rz_re_r\mapsto\sum_rz_rr\) is integral linear;
it sends the present vector \(e_0\) to zero.  These statements require no
map from \(\C\) or \(\mathbb Q\) to the finite group.

\section{Original masses and an evaluated determinant loss}

For the metric statements extend the preceding groups and maps to \(\C\).
Place the orthonormal metric on the original word basis of \(\C[S]\).
Define the positive masses only on their actual images:
\[
 \mu_S(x)=|\{c\in S:v(c)=x\}|,\qquad
 \bar\mu_S(r)=\sum_{\substack{x\in X_S\\x\bmod b=r}}\mu_S(x).
 \tag{E9}
\]
No zero-mass value is inserted into an inverse diagonal matrix.

\begin{theorem}[Exact integral-section loss]
The quotient metric on \(E_S\otimes\C\), the residue quotient metric,
and the unique minimum-norm right inverse of \(\pi_S\) are respectively
\[
 G_{E,S}=\operatorname{diag}_{x\in X_S}\frac1{\mu_S(x)},\quad
 G_{\bar E,S}=\operatorname{diag}_{r\in R_S}\frac1{\bar\mu_S(r)},\quad
 s_S(e_r)=\sum_{x\bmod b=r}\frac{\mu_S(x)}{\bar\mu_S(r)}e_x .
 \tag{E10}
\]
Let \(d_S(r)\) be the least actual integer value in the residue fibre
and let \(\sigma_S(e_r)=e_{d_S(r)}\).  Put \(K_S=\sigma_S-s_S\).
Then
\[
 \begin{gathered}
 \pi_SK_S=0,\qquad
 \sigma_S^*G_{E,S}\sigma_S
   =G_{\bar E,S}+K_S^*G_{E,S}K_S,\\
 K_S^*G_{E,S}K_S
   =\operatorname{diag}_{r\in R_S}
       \left(\frac1{\mu_S(d_S(r))}-\frac1{\bar\mu_S(r)}\right),\\
 \log\frac{\det(\sigma_S^*G_{E,S}\sigma_S)}{\det G_{\bar E,S}}
   =\sum_{r\in R_S}\log\frac{\bar\mu_S(r)}{\mu_S(d_S(r))}.
 \end{gathered}
 \tag{E11}
\]
Every determinant is in the displayed original residue basis.  Empty
determinants are one.
\end{theorem}
\begin{proof}
For a fixed \(x\), minimizing \(\sum_{v(c)=x}|z_c|^2\) subject to
\(\sum_{v(c)=x}z_c=u_x\) gives the unique solution
\(z_c=u_x/\mu_S(x)\) and cost \(|u_x|^2/\mu_S(x)\), by completing
the sum of squared deviations.  Different value fibres are disjoint.
This proves \(G_{E,S}\).

Within residue \(r\), write \(u_x=\mu_S(x)t/\bar\mu_S(r)+h_x\)
where \(\sum_{x\bmod b=r}h_x=0\).  Direct expansion gives
\[
 \sum_{x\bmod b=r}\frac{|u_x|^2}{\mu_S(x)}
 =\frac{|t|^2}{\bar\mu_S(r)}
  +\sum_{x\bmod b=r}\frac{|h_x|^2}{\mu_S(x)} .
 \tag{E12}
\]
The cross term is zero because its coefficient is
\(\sum_xh_x=0\).  The displayed cost is uniquely minimized by \(h=0\).
This proves the last two parts of (E10), as well as orthogonality of
\(\operatorname{im}s_S\) to \(\ker\pi_S\).  Both \(s_S\) and
\(\sigma_S\) are right inverses, so \(K_S\) has image in that kernel.
Substituting \(h=K_St\) in (E12) proves the first identity in (E11).
Columns for different residues have disjoint value support.
The squared norm of \(\sigma_Se_r\) is
\(1/\mu_S(d_S(r))\), while that of \(s_Se_r\) is
\(1/\bar\mu_S(r)\).  This proves the diagonal formula, and taking
the quotient of the diagonal determinants proves the final formula.
\end{proof}

Thus the squared generalized singular values of the actual integral
section relative to the attained quotient metric are
\(\bar\mu_S(r)/\mu_S(d_S(r))\).  Formula (E11) evaluates its full loss,
including every original representation multiplicity.

\section{Source enlargement has an exact metric boundary}

Write \(i_{ST}:E_S\otimes\C\to E_T\otimes\C\) and
\(\bar i_{ST}:\bar E_S\otimes\C\to\bar E_T\otimes\C\) for the original
basis inclusions when \(S\subseteq T\).  Extend \(\mu_S(x)\) by zero
for \(x\in X_T\setminus X_S\).  The minimum-norm section defect is the
map with its stated source and target
\[
 K_{ST}=i_{ST}s_S-s_T\bar i_{ST}:
    \bar E_S\otimes\C\longrightarrow\ker\pi_T.
 \tag{E13}
\]
\begin{proposition}[Full source-change norm and its composition law]
For every \(r\in R_S\),
\[
 \begin{split}
 K_{ST}e_r
   &=\sum_{\substack{x\in X_T\\x\bmod b=r}}
      \left(\frac{\mu_S(x)}{\bar\mu_S(r)}
           -\frac{\mu_T(x)}{\bar\mu_T(r)}\right)e_x,\\
 \|K_{ST}e_r\|_{G_{E,T}}^2
   &=\frac{1}{\bar\mu_S(r)^2}
        \sum_{x\bmod b=r}\frac{\mu_S(x)^2}{\mu_T(x)}
        -\frac1{\bar\mu_T(r)},\\
 0\leq\|K_{ST}e_r\|_{G_{E,T}}^2
   &\leq\frac1{\bar\mu_S(r)}-\frac1{\bar\mu_T(r)}.
 \end{split}
 \tag{E14}
\]
The columns for distinct residues are orthogonal in \(G_{E,T}\).  For
\(S\subseteq T\subseteq U\) the exact typed composition identity is
\[
 K_{SU}=i_{TU}K_{ST}+K_{TU}\bar i_{ST}.
 \tag{E15}
\]
\end{proposition}
\begin{proof}
The two right-inverse identities and the commutative square give
\(\pi_TK_{ST}=\bar i_{ST}-\bar i_{ST}=0\).
Insertion of (E10) gives the first formula.  In the \(G_{E,T}\) norm,
expansion of its square has cross term
\(-2/\bar\mu_T(r)\) and last term \(1/\bar\mu_T(r)\), yielding the
second formula.  Nonnegativity follows either from this norm or
Cauchy--Schwarz.  Since \(S\subseteq T\) gives
\(0\leq\mu_S(x)\leq\mu_T(x)\), one has
\(\sum_x\mu_S(x)^2/\mu_T(x)\leq\bar\mu_S(r)\), giving the upper bound.
The disjoint residue supports prove column orthogonality.  Finally,
\[
 i_{TU}(i_{ST}s_S-s_T\bar i_{ST})
 +(i_{TU}s_T-s_U\bar i_{TU})\bar i_{ST}
 =i_{SU}s_S-s_U\bar i_{SU},
 \]
which is (E15), with all domains as in (E13).
\end{proof}

The integral sections have the parallel defect
\[
 i_{ST}\sigma_Se_r-\sigma_T\bar i_{ST}e_r
     =e_{d_S(r)}-e_{d_T(r)}\in\ker\pi_T .
 \tag{E16}
\]
Its squared \(G_{E,T}\)-norm is zero if the representatives coincide,
and otherwise is \(1/\mu_T(d_S(r))+1/\mu_T(d_T(r))\), because those
are two different original basis vectors.  The same addition and
subtraction proves its version of (E15).

\section{A universal derivation of the seven arithmetic matrices}

The composition source retains
\[
 A=(10,\{1,3\}),\quad B=(10,\{2,4\}),\quad
 D_A=H_5(\{1,3\})=[0,16],\quad
 D_B=H_5(\{2,4\})=2[0,12].
 \tag{E17}
\]
For \(A\), the intervals \([3j,3j+4]\), \(0\leq j\leq4\), overlap
and fill \([0,16]\).  For \(B\), the original values are
\(2(i+2j)\) with \(0\leq i,j\leq4\); the intervals
\([2j,2j+4]\) overlap and fill \([0,12]\).  Thus (E17) is an equality
of original numerical sets.

Keep the seven source states in their stated order:
\[
 \mathcal R=((0),(0,1),(0,2),(0,1,2),(0,3),(0,1,3),(0,1,2,3)).
 \tag{E18}
\]
For any finite \(Y\subset\Z\), put \(f_R(Y)=|Y+R|\).  For
\(X=A,B\), \(R\in\mathcal R\), and \(j\in\{0,\ldots,9\}\), define
\[
 C_{X,R,j}
  =\left\{\frac{d+r-j}{10}:
        d\in D_X,\ r\in R,\ d+r\equiv j\pmod {10}\right\}.
 \tag{E19}
\]
If this set is nonempty, write \(m_{X,R,j}=\min C_{X,R,j}\) and
\(R_{X,R,j}=C_{X,R,j}-m_{X,R,j}\).
The change of presentation is the explicit bijection
\[
 Y+R_{X,R,j}\longrightarrow
 (D_X+10Y+R)\cap(j+10\Z),\qquad
 z\longmapsto j+10(m_{X,R,j}+z),
 \tag{E20}
\]
whose inverse is \(t\mapsto(t-j)/10-m_{X,R,j}\).  Each shift and
residue is retained in \texttt{UNIVERSAL\_INTERFACE\_CHECKS.json}.
Direct insertion of the finite sets in (E17)--(E19) gives the following
numbers of residues of each resulting state:
\[
 M_A=\begin{pmatrix}
3&7&0&0&0&0&0\\2&8&0&0&0&0&0\\1&9&0&0&0&0&0\\
1&9&0&0&0&0&0\\0&10&0&0&0&0&0\\0&10&0&0&0&0&0\\
0&10&0&0&0&0&0
\end{pmatrix},
\quad
 M_B=\begin{pmatrix}
0&2&0&3&0&0&0\\0&4&0&6&0&0&0\\0&1&0&4&0&0&0\\
0&3&0&7&0&0&0\\0&4&0&6&0&0&0\\0&3&0&7&0&0&0\\
0&2&0&8&0&0&0
\end{pmatrix}.
 \tag{E21}
\]
For completeness these entries can be obtained without a numerical tail:
enumerate each of the at most \(25\cdot4\) displayed pairs \((d,r)\),
place \((d+r-j)/10\) into its residue set in (E19), subtract its recorded
minimum with inverse (E20), and count the state.  The accompanying
independent verifier records all \(140\) such original fibres and
checks each entry of (E21).

The residues in (E20) are disjoint.  Summing their cardinalities proves,
for \emph{every} finite numerical tail \(Y\), the identities
\[
 f(D_X+10Y)=M_Xf(Y),\qquad
 f(\{0\})=(1,2,2,3,2,3,4)^{\mathsf T}.
 \tag{E22}
\]
The matrices thus describe the exact count observations of the retained
numerical sets.  They are induced through the original residue partition
and bijections (E20); no change to coefficient-word amplitudes is involved.
Multiplication of the displayed integer matrices gives
\[
 M_AM_B=(17,18,19,19,20,20,20)^{\mathsf T}(0,2,0,3,0,0,0).
 \tag{E23}
\]
This supplies the universal finite-source input to the supplied
composition proofs, alongside their exact finite and limiting attainers.

\section{The retained empty radix and its actual kernel}

Retaining the two \(A\)-positions of \(ABA\) gives
\(\{1,3,100,300\}\).  The distinct local values are
\([0,16]^2\), and the two evaluations at issue are
\[
 e_{x,y}\longmapsto e_{x+100y},\qquad
 e_{x,y}\longmapsto e_{x+10y}.
 \tag{E24}
\]
The first is a bijection onto a \(289\)-element basis: if its values
coincide, \(100(y-y')=x'-x\in[-16,16]\), forcing both differences
to vanish.  The second has image \([0,176]\), because the intervals
\([10y,10y+16]\), \(0\leq y\leq16\), overlap and have precisely these
endpoints.  It is therefore an onto map from a free rank-\(289\) group
to a free rank-\(177\) group, with free kernel rank \(112\).
An explicit kernel basis chooses for each value \(z\) one pair
\((x_z,y_z)\) evaluating to \(z\) and uses
\(e_{x,y}-e_{x_z,y_z}\) for every other pair in that fibre.
Within each fibre these vectors are independent and span its
coefficient-sum-zero group, proving the rank calculation integrally.

The first evaluation identifies its basis with the original retained
image, so composing its inverse with the second evaluation gives the
exact contraction map between the two value modules.  Formula (E6)
then gives its corresponding additional receiving kernel.  The mass of
each local value pair remains the product of the two original local
word multiplicities.  Applying (E10)--(E12) to the fibres of (E24)
supplies its original minimum-norm quotient, with no unit-mass
substitution.

\section{Accepted return and next arithmetic scope}

The preserved paper proves sharp ternary residue prices
\(3,5,13,23\) for ranks \(1,2,3,4\) of individually modular-five- or
modular-six-free canonical blocks, at unrestricted height and radix.
Its classical addition-theorem input is Martin Kneser's stabilizer
theorem, in the stated form attributed to Matt DeVos's primary proof
\cite{devos}.  The remaining \(490\) exceptional source blocks have
explicit original subset-mask witnesses, replayed in this intake.
The preserved composition paper also gives every fixed-composition
minimum and the full limiting frequency spectrum in the actual
radix-ten dictionary.  All supplied credit to the original
anonymous/deleted contributor and to \texttt{sneed-and-feed}'s distinct
Lean extension remains in the original paper.

The present note proves the exact finite cochain realization (E6),
the integral-section determinant loss (E11), the source-change
boundary and composition law (E14)--(E15), and the universal original
matrix return (E19)--(E23).  These are usable instances of the retained
Split-Zero supported quotient, with their full metrics specified.
The larger-rank modular arithmetic price remains an uncomputed part
of this EP817 continuation.  The original theta source and its signed
analytic allocations receive the algebraic construction (E6)--(E16);
no matrix identifying those analytic sources with the finite residue
source has been supplied or proved in this intake.  Therefore no
analytic allocation coefficient or statement about zeta zeros is
entered in the accepted-result index from this package.

\begin{thebibliography}{9}
\bibitem{continuation}
The Clankers, \emph{Residue prices, constrained composition, and nonperiodic
optimal schedules}, 17 September 2026.  User-supplied
\texttt{EP817\_CONTINUATION\_20260917}; full source preserved at
\texttt{source/research/residue\_composition/}.
Baseline repository
\url{https://github.com/KokunoYumeto/erdos-problem-817-workbench},
revision \texttt{dbd0a93f2cf935103e88e5c2b2b71fe85ae7537b}.
\bibitem{support}
KokunoYumeto, \emph{Reconstruction with changes of support index},
\texttt{workbenches/splitzero-tandem/tex/support\_diagrams.tex},
equations (D6)--(D8), revision
\texttt{58626a62cd5648e8ae7450fb54d4a4aab2330981},
\url{https://github.com/KokunoYumeto/zeta-function-research-reader}.
The pinned source and the retained local version with (D9)--(D14) were
read during this integration.
\bibitem{devos}
Matt DeVos, \emph{A short proof of Kneser's addition theorem for abelian
groups}, 2013, \url{https://arxiv.org/abs/1303.3539v1}.
Role: the classical stabilizer inequality in the preserved residue-price
proof; the primary statement was checked on 19 September 2026.
\end{thebibliography}
\end{document}
